\subsection{COMPOSITION OPERATORS ON AN ALGEBRA OF POLYNOMIALS}

Let K be a commutative field of characteristic 0, and K{[}X{]} the algebra of polynomials in one indeterminate over K (Alg., IV. 1); throughout this section by an operator on K{[}X{]} we shall mean a linear map U of the vector space K{[}X{]} (over K) into itself; since the monomials \(X^{n}\) ( \(n \geqslant 0\) ) form a basis for this space, U is determined by the polynomials \(U(X^{n})\) ; specifically, if \(f(X) = \sum_{k=0}^{\infty} \lambda_{k} X^{k}\) with \(\lambda_{k} \in K\) , then \(U(f) = \sum_{k=0}^{\infty} \lambda_{k} U(X^{k})\) .

If \(\mathbf{G}\) is a commutative algebra over \(\mathbf{K}\) , with an identity element, the \(\mathbf{G}\) -module \(\mathrm{G}[\mathrm{X}]\) is obtained by extending the field of scalars \(\mathbf{K}\) of the vector space \(\mathrm{K}[\mathrm{X}]\) to \(\mathrm{G}\) ; every operator \(U\) on \(\mathrm{K}[\mathrm{X}]\) extends in a unique manner to a linear map of the \(\mathrm{G}\) -module \(\mathrm{G}[\mathrm{X}]\) to itself, which we again denote by \(U\) (Alg., II, p.~278); for each element \(g(\mathrm{X}) = \sum_{k=0}^{\infty} \gamma_k X^k\) , with \(\gamma_k \in \mathrm{G}\) , one has \(U(g) = \sum_{k=0}^{\infty} \gamma_k U(X^k)\) .

Consider in particular the case where \(\mathrm{G} = \mathrm{K}[\mathrm{Y}]\) ; then \(\mathrm{G}[\mathrm{X}]\) is the ring \(\mathrm{K}[\mathrm{X},\mathrm{Y}]\) of polynomials in two indeterminates over \(\mathrm{K}\) ; to avoid confusion we denote the extension of \(U\) to \(\mathrm{G}[\mathrm{X}]\) by \(U_{\mathrm{X}}\) . For any polynomial \(g(\mathrm{X},\mathrm{Y}) = \sum_{k=0}^{\infty}\gamma_k(\mathrm{Y})\mathrm{X}^k\) where \(\gamma_k(Y)\in \mathrm{K}[\mathrm{Y}]\) one thus has \(U_{\mathrm{X}}(g) = \sum_{k=0}^{\infty}\gamma_k(\mathrm{Y})U(\mathrm{X}^k)\) . Since \(U_{\mathrm{X}}\) is linear one sees that if one writes \(g(\mathrm{X},\mathrm{Y}) = \sum_{h=0}^{\infty}\beta_h(\mathrm{X})\mathrm{Y}^h\) then one also has \(U_{\mathrm{X}}(g) = \sum_{h=0}^{\infty}U(\beta_h)\mathrm{Y}^h\) .

Under the canonical isomorphism of K{[}X{]} onto K{[}Y{]} which associates Y with X the operator U transforms to an operator on K{[}Y{]}, which we will denote \(U_{Y}\) , to avoid confusion, so that \(U_{\mathrm{Y}}(f(\mathrm{Y}))\) is the polynomial obtained by replacing X by Y in the polynomial \(U(f(\mathrm{X})) = U_{\mathrm{X}}(f(\mathrm{X}))\) . This operator \(U_{Y}\) can in turn be extended to an operator (again denoted by \(U_{Y}\) ) on K{[}X,Y{]}: if \(g(\mathrm{X}, \mathrm{Y}) = \sum_{h=0}^{\infty} \beta_{h}(\mathrm{X}) \mathrm{Y}^{h}\) then \(U_{\mathrm{Y}}(g(\mathrm{X}, \mathrm{Y})) = \sum_{h=0}^{\infty} \beta_{h}(\mathrm{X}) U_{\mathrm{Y}}(\mathrm{Y}^{h})\) .

As an example of these extensions we cite the derivation operator D on K{[}X{]} (Alg., IV. 6), which gives the partial derivation operators \(D_{X}\) and \(D_{Y}\) on K{[}X,Y{]}.

For any polynomial \(f \in K[X]\) we denote by \(T_{Y}(f)\) the polynomial \(f(X + Y)\) in K{[}X,Y{]}; the map \(T_{Y}\) is a K-linear map from K{[}X{]} into K{[}X,Y{]}, called a translation operator.

DEFINITION 1. One says that an operator U on K{[}X{]} is a composition operator if it commutes with the translation operator, that is, if \(U_{X}T_{Y} = T_{Y}U\) .

In other words, if \(f\) is an arbitrary polynomial in \(\mathbf{K}[\mathbf{X}]\) , and if \(g = U(f)\) , one must have \(g(\mathbf{X} + \mathbf{Y}) = U_{\mathbf{X}}(f(\mathbf{X} + \mathbf{Y}))\) .

It follows immediately from this definition that for every polynomial \(f(\mathrm{X}) \in \mathrm{K}[\mathrm{X}]\) one has, in the above notation,

\[
U _ {\mathrm{X}} \big (f (\mathrm{X} + \mathrm{Y}) \big) = U _ {\mathrm{Y}} \big (f (\mathrm{X} + \mathrm{Y}) \big).\tag{1}
\]

Examples. 1) For every \(\lambda \in \mathbf{K}\) the operator which to every polynomial \(f(\mathbf{X})\) associates the polynomial \(f(\mathbf{X} + \lambda)\) is a composition operator.

\begin{enumerate}
\def\labelenumi{\arabic{enumi})}
\setcounter{enumi}{1}
\tightlist
\item
  The derivation D on K{[}X{]} is a composition operator (cf.~prop. 1).
\end{enumerate}

Remark. Since K is an infinite field the operator U on K{[}X{]} is a composition operator if and only if for every polynomial \(f \in K[X]\) and every element \(\alpha \in K\) one has, putting \(g = U(f)\) , that \(g(X + \alpha) = U(f(X + \alpha))\) . (Alg., IV. 16, cor.).

It is clear that every linear combination of composition operators, with coefficients in K, is a composition operator; the same is true for the composition of two composition operators. In other words, the composition operators form a subalgebra \(\Gamma\) of the algebra of endomorphisms of the vector space K{[}X{]}.

PROPOSITION 1. For an operator U on K{[}X{]} to be a composition operator it is necessary and sufficient that it commute with the derivation D on K{[}X{]}.

Indeed, the Taylor formula shows that for every polynomial \(f \in K[X]\) one has

\[
U _ {\mathrm{X}} \big (f (\mathrm{X} + \mathrm{Y}) \big) = U _ {\mathrm{X}} \left(\sum_ {k = 0} ^ {\infty} \frac {1}{k !} \mathrm{Y} ^ {k} \mathrm{D} ^ {k} f (\mathrm{X})\right) = \sum_ {k = 0} ^ {\infty} \frac {1}{k !} \mathrm{Y} ^ {k} U \big (\mathrm{D} ^ {k} f (\mathrm{X}) \big);
\]

if one puts \(g = U(f)\) then one has

\[
g (\mathrm{X} + \mathrm{Y}) = \sum_ {k = 0} ^ {\infty} \frac {1}{k !} \mathrm{Y} ^ {k} \mathrm{D} ^ {k} g (\mathrm{X}) = \sum_ {k = 0} ^ {\infty} \frac {1}{k !} \mathrm{Y} ^ {k} \mathrm{D} ^ {k} (U (f (\mathrm{X})))
\]

for U to be a composition operator, one must then have \(UD^{k} = D^{k}U\) for every integer \(k \geqslant 1\) , and in particular \(UD = DU\) . Conversely, if this relation holds, it implies \(UD^{k} = D^{k}U\) for every integer \(k \geqslant 1\) , by induction on k; the Taylor formula then shows that \(g(X + Y) = U_{X}(f(X + Y))\) .

For every polynomial \(f \in K[X, Y]\) we denote by \(U_{0}(f)\) the term independent of X in the polynomial \(U_{X}(f)\) ; in particular, if \(f \in K[X]\) , \(U_{0}(f)\) is the constant term in \(U(f)\) , and \(U_{0}\) is a linear form on K{[}X{]}. For every polynomial \(f \in K[X]\) let \(g = U(f)\) ; by def. 1 of VI, p.~270

\[
g (\mathrm{X} + \mathrm{Y}) = U _ {\mathrm{X}} (f (\mathrm{X} + \mathrm{Y})) = U _ {\mathrm{X}} \left(\sum_ {k = 0} ^ {\infty} \frac {1}{k !} \mathrm{X} ^ {k} \mathrm{D} ^ {k} f (\mathrm{Y})\right) = \sum_ {k = 0} ^ {\infty} \frac {1}{k !} U (\mathrm{X} ^ {k}) \mathrm{D} ^ {k} f (\mathrm{Y})
\]

and if, in this formula, one replaces X by 0, one obtains

\[
g (\mathrm{Y}) = \sum_ {k = 0} ^ {\infty} \frac {1}{k !} U _ {0} (\mathrm{X} ^ {k}) \mathrm{D} ^ {k} f (\mathrm{Y}).
\]

Thus one sees that

\[
U (f (\mathrm{X})) = \sum_ {k = 0} ^ {\infty} \frac {1}{k !} \mu_ {k} \mathrm{D} ^ {k} f (\mathrm{X})\tag{2}
\]

where \(\mu_{k}\) is the constant term in the polynomial \(U(\mathbf{X}^k)\) .

This formula shows that the \(\mu_{k}\) determine the composition operator U completely; conversely, if \((\mu_{n})\) is an arbitrary sequence of elements of K then the formula (2) defines an operator U which clearly commutes with D, so (VI, p.~270, prop. 1) is a composition operator. From now on we shall write (2) in the form

\[
U = \sum_ {k = 0} ^ {\infty} \frac {1}{k !} \mu_ {k} D ^ {k}.\tag{3}
\]

This formula can be interpreted in topological language as follows: if one considers the discrete topology on K{[}X{]}, and the topology of simple convergence on the algebra \(\operatorname{End}(\operatorname{K}[X])\) of endomorphisms of K{[}X{]} (Gen.~Top., X, p.~277), the series with general term \(\frac{1}{k!}\mu_{k}D^{k}\) is commutatively convergent in \(\operatorname{End}(\operatorname{K}[X])\) and has sum U (Gen.~Top., III, p.~269).

The formula (3) shows that to every formal series \(u(S) = \sum_{k=0}^{\infty} \alpha_k S^k\) in one indeterminate over \(K\) (Alg., IV. 41) one can associate the composition operator \(U = \sum_{k=0}^{\infty} \alpha_k D^k\) , which in future we shall denote by \(u(D)\) . This remark can be clarified in the following manner:

THEOREM 1. The map which to every formal series \(u(\mathbf{S}) = \sum_{k=0}^{\infty} \alpha_k \mathbf{S}^k\) in one indeterminate over \(\mathbf{K}\) associates the composition operator \(u(\mathbf{D}) = \sum_{k=0}^{\infty} \alpha_k \mathbf{D}^k\) on \(\mathrm{K}[\mathrm{X}]\) is an isomorphism of the algebra \(\mathrm{K}[[\mathrm{S}]]\) of formal series onto the algebra \(\Gamma\) of composition operators.

One verifies immediately that this map is an homomorphism. It remains to see that it is injective, in other words, that the relation \(\sum_{k=0}^{\infty}\alpha_{k}D^{k}=0\) implies \(\alpha_{k}=0\) for every k; now \(h!\alpha_{h}\) is the constant term in the polynomial obtained by applying \(\sum_{k=0}^{\infty}\alpha_{k}D^{k}\) to \(X^{h}\) , whence the theorem.

COROLLARY. The algebra \(\Gamma\) of composition operators in K{[}X{]} is commutative.

Example. If \(U\) is the operator which to each polynomial \(f(\mathrm{X})\) associates \(f(\mathrm{X} + \lambda)\) (where \(\lambda \in \mathbf{K}\) ), one has \(U_0(\mathrm{X}^k) = \lambda^k\) , and so \(U = \sum_{k=0}^{\infty} \frac{1}{k!} (\lambda \mathrm{D})^k\) . By analogy with the series expansion of \(e^x\) (III, p.~105) we write \(e^{\mathrm{S}}\) or \(\exp(\mathrm{S})\) for the formal series \(\sum_{n=0}^{\infty} \frac{1}{n!} \mathrm{S}^n\) in the ring \(\mathrm{K}[[\mathrm{S}]]\) ; so one can write \(U = e^{\lambda \mathrm{D}}\) . Replacing the field \(\mathrm{K}\) by the field of rational fractions \(\mathrm{K}(\mathrm{Y})\) in this argument one sees similarly that the translation operator \(T_{\mathrm{Y}}\) can be written \(e^{\mathrm{YD}}\) .

Furthermore, in the ring K{[}{[}S, T{]}{]} of formal series in two indeterminates over K one has

\[
\begin{array}{l} (\exp S) (\exp T) = \sum_ {p, q} \frac {S ^ {p}}{p !} \frac {T ^ {q}}{q !} \\ = \sum_ {n = 0} ^ {\infty} \frac {1}{n !} \left(\binom {n} {0} S ^ {n} + \binom {n} {1} S ^ {n - 1} T + \dots + \binom {n} {n} T ^ {n}\right) \\ = \exp (S + T) \end{array}\tag{4}
\]

and in particular

\[
(\exp S) (\exp (- S)) = 1\tag{5}
\]

which justifies the notation we have introduced.

Scholium. The isomorphism of the algebra K{[}{[}S{]}{]} of formal series and the algebra \(\Gamma\) of composition operators on K{[}X{]} sometimes allows one to prove propositions on formal series most simply, by proving them for the corresponding composition operators (cf.~VI, p.~274, prop. 6).

